\section{Referencial Teórico}
\label{referencial}

\textbf{Grafos e redes.} Um grafo ponderado é $G=(V,E,w)$, com vértices $V$, arestas $E \subseteq V\times V$ e uma função $w: E \rightarrow \mathbb{R}^+$ que expressa a intensidade da relação \citep{barabasi2016}. Em redes sociais, vértices representam indivíduos e arestas representam relações; modelos \textit{multirrelacionais} combinam diferentes tipos de relação com pesos distintos \citep{porto2025}. Quando a relação é de \textit{similaridade} entre registros tabulares, é comum conectar cada vértice aos seus $k$ vizinhos mais similares (grafo $k$-NN), o que produz uma rede esparsa e conexa. Para atributos mistos (ordinais e nominais), o coeficiente de \citet{gower1971} agrega similaridades por atributo em $[0,1]$.

\textbf{Homofilia e assortatividade.} Homofilia é a tendência de vértices semelhantes se conectarem \citep{mcpherson2001}. O coeficiente de assortatividade por atributo categórico \citep{newman2003} é
$r = (\sum_i e_{ii} - \sum_i a_i b_i)/(1-\sum_i a_i b_i)$,
em que $e_{ij}$ é a fração de arestas entre as categorias $i$ e $j$ e $a_i$, $b_i$ são as frações de extremidades de arestas na categoria $i$. Tem-se $r=0$ para mistura aleatória e $r=1$ para segregação perfeita.

\textbf{Comunidades.} Uma comunidade é um subgrafo com densidade interna significativamente maior que a densidade com o restante da rede \citep{fortunato2010}. A qualidade de uma partição é usualmente medida pela modularidade \citep{newman2004}:
$Q=\frac{1}{2m}\sum_{ij}\left[A_{ij}-\frac{k_ik_j}{2m}\right]\delta(c_i,c_j)$.
Neste trabalho comparam-se três métodos: \textit{Louvain/Multilevel} \citep{blondel2008}, que otimiza $Q$ por movimentos locais e agregação hierárquica em $O(|V|\log|V|)$; \textit{Fast Greedy} \citep{clauset2004}, que funde comunidades de forma gulosa em $O(|V|\log^2|V|)$; e \textit{Label Propagation} \citep{raghavan2007}, que propaga rótulos por maioria de vizinhos em tempo quase linear.

\textbf{Aprendizado sobre grafos.} Em \textit{propagação de rótulos} semissupervisionada \citep{zhou2004}, os rótulos conhecidos difundem-se pelas arestas segundo $F^{(t+1)}=\alpha S F^{(t)}+(1-\alpha)Y$, com $S=D^{-1/2}WD^{-1/2}$, de modo que vértices densamente conectados recebem rótulos semelhantes. \textbf{Lift} mede a associação entre dois eventos, $\mathrm{lift}(a,b)=P(a\wedge b)/(P(a)P(b))$, sendo $\mathrm{lift}>1$ indicativo de co-ocorrência acima do acaso \citep{agrawal1993}.

\textbf{Trabalhos relacionados.} \citet{tinto1975} consolidou o modelo de integração acadêmica e social da evasão, que enfatiza a interação entre atributos de entrada e a experiência no curso. \citet{realinho2022} publicaram o dataset aqui utilizado e mostraram que o desempenho nos primeiros semestres é o principal preditor do desfecho. \citet{romero2020} revisam a Mineração de Dados Educacionais e destacam a predição de evasão como uma de suas tarefas centrais, majoritariamente abordada com classificadores tabulares. \citet{porto2025} propuseram três modelos de grafo de similaridade (contagem direta, multirrelacional e medida ordinal) para candidatos a assistência estudantil e aplicaram Fast Greedy, Multilevel e Walktrap; o modelo multirrelacional obteve as comunidades mais coesas. \citet{acuna2025} propuseram o CRISP-NET, que integra as fases do CRISP-DM à análise de redes. Este trabalho adapta o modelo multirrelacional de \citet{porto2025} ao problema de evasão, acrescenta uma rede de fatores e valida a estrutura contra o desfecho observado.
